\begin{lemma}[Stability of a color swap outside a closed tree]
Let \(T\) be a tree-sequence, let \(\varphi\) be a proper partial
\(k\)-edge-coloring leaving exactly \(e_0\) uncolored, and suppose that
\(V(T)\) is closed under \(\varphi\).  Let \(\alpha\ne\beta\) be two
colors missing on \(T\), and let \(C\subseteq[k]\) be arbitrary.  Define
\[
 \psi=\noderef{OutsideTreeColorSwap}(T,\varphi,\alpha,\beta).
\]
Then \(\psi\) is proper, leaves exactly \(e_0\) uncolored, and is
\((T,C,\varphi)\)-stable.  Moreover, every exact edge identity incident
with \(T\) has the same color under \(\psi\) and \(\varphi\).
\end{lemma}
\begin{proof}
Write \(X=V(T)\).  We first show that no boundary edge of \(X\) has color
\(\alpha\) or \(\beta\) under \(\varphi\).  By
\(\noderef{TreeMissingColors}\), choose \(x_\alpha,x_\beta\in X\) at
which \(\alpha,\beta\), respectively, are missing.  If an exact boundary
identity \(f\) had, say, color \(\alpha\), then the closedness condition
\(\noderef{ColoredVertexSetClosed}\), applied to \(f,x_\alpha\), would
say that \(\alpha\) is not missing at \(x_\alpha\), a contradiction.
The same argument applies to \(\beta\).  Hence
\[
 f\in\noderef{BoundaryEdges}(G,X)
 \quad\Longrightarrow\quad
 \varphi(f)\notin\{\mathsf{some}(\alpha),
                       \mathsf{some}(\beta)\}.                 \tag{1}
\]
By the definition of \(\noderef{OutsideTreeColorSwap}\), every
\(\noderef{TreeIncidentEdge}\)-identity is fixed.  On every other edge,
the same transposition \(\alpha\leftrightarrow\beta\) is applied, and
all other values, including \(\mathsf{none}\), are fixed.  In particular,
\[
 \psi(f)=\mathsf{none}\quad\Longleftrightarrow\quad
 \varphi(f)=\mathsf{none}.                                   \tag{2}
\]
We verify properness using exact edge identities.  Let \(f\ne g\) be
adjacent.  If both are incident with \(T\), both are fixed, so equality
of their new nonempty colors would contradict the properness of
\(\varphi\) in \(\noderef{IsProperEdgeColoring}\).  If neither is
incident with \(T\), both old values undergo the same transposition.
That transposition is injective because \(\alpha\ne\beta\): it exchanges
the two displayed values and fixes every other option value.  Equality of
the new colors would therefore again give equality of the old colors,
contrary to properness.

It remains to consider the mixed case.  Suppose that \(f\) is incident
with \(T\) and \(g\) is not.  By \(\noderef{AdjacentEdges}\), they share
a vertex \(w\).  The vertex \(w\) is outside \(X\), since otherwise \(g\)
would be incident with \(T\).  Choose \(x\in X\) incident with \(f\).
Then \(x\ne w\), and \(\noderef{Incident}\) together with the two
endpoint description in \(\noderef{LooplessMultigraph}\) shows that its
endpoint pair is exactly \(\{x,w\}\); hence
\(\noderef{BoundaryEdges}\) shows that \(f\) is
a boundary identity of \(X\).  Consequently (1) says that the old and
new color of \(f\) is neither \(\alpha\) nor \(\beta\).  The outside
transposition can give \(g\) this same other color only if \(g\) already
had it.  Thus a new equality on \(f,g\) would have been an old equality,
again contradicting properness.  Interchanging \(f\) and \(g\) handles
the other mixed orientation.  This proves that \(\psi\) is proper.

Equation (2) and the two clauses of
\(\noderef{IsOneEdgeUncolored}\) show that precisely the same identity
\(e_0\) is uncolored under \(\psi\).

Finally, every edge incident with a vertex \(x\in X\) is a tree-incident
edge and is therefore fixed.  The incidence-witness definition of
\(\noderef{PresentColors}\) gives
\(\operatorname{present}_\psi(x)=\operatorname{present}_\varphi(x)\),
and taking complements using \(\noderef{MissingColors}\) gives
\[
                  \overline\psi(x)=\overline\varphi(x)
                  \qquad(x\in X).                            \tag{3}
\]
The first clause of \(\noderef{StableColoring}\) holds because every
tree-incident edge is fixed, without any restriction on its color; the
second clause is exactly (3).  Thus \(\psi\) is
\((T,C,\varphi)\)-stable.  The asserted pointwise equality on all
tree-incident identities was already obtained directly from the outside
swap definition, completing the proof.
\end{proof}
